\section{Dữ liệu và phương pháp đánh giá}

\subsection{Dữ liệu}

\subsubsection{Nguồn dữ liệu}

Bộ dữ liệu ban đầu của nghiên cứu là NewsQA, được tải từ kho
\href{https://huggingface.co/datasets/lucadiliello/newsqa}
{\texttt{lucadiliello/newsqa}} trên Hugging Face. NewsQA gồm các bài báo CNN,
câu hỏi do người đọc đặt, đáp án dạng đoạn văn bản ngắn (answer span) và vị trí
ký tự của đoạn đáp án trong bài báo.

Dữ liệu được tổ chức theo câu hỏi, không gom sẵn theo bài báo. Mỗi dòng chứa
một câu hỏi và toàn bộ bài báo tương ứng. Vì vậy, một bài có nhiều câu hỏi sẽ
được lặp ở nhiều dòng. Nguồn không có \texttt{article\_id}; số bài được tính
bằng số \texttt{context} khác nhau.

Ví dụ rút gọn một dòng dữ liệu gốc:

\begin{verbatim}
{
  "context": "... full CNN article ...",
  "question": "When was Pandher sentenced to death?",
  "answers": ["February."],
  "key": "724f6eb9a2814e4fb2d7d8e4de846073",
  "labels": [{"start": [261], "end": [269]}]
}
\end{verbatim}

Ý nghĩa các trường:

\begin{itemize}
  \item \texttt{context}: toàn bộ nội dung bài báo CNN.
  \item \texttt{question}: câu hỏi được đặt dựa trên bài báo đó.
  \item \texttt{answers}: một hoặc nhiều đáp án ngắn lấy từ bài báo.
  \item \texttt{key}: mã định danh của câu hỏi.
  \item \texttt{labels}: vị trí ký tự đầu và cuối của đáp án trong
  \texttt{context}. Ví dụ, \texttt{start = 261} và \texttt{end = 269} nghĩa
  là đáp án gồm cả ký tự ở vị trí 269. Pipeline đổi thành đoạn
  $[261,270)$ và kiểm tra đoạn này khớp với đáp án.
\end{itemize}

Revision nguồn được khóa tại
\texttt{728e52920b8e4ffcfaad93fa47556f26a1d82546}. Thống kê theo định dạng gốc
như sau:

\begin{table}[htbp]
  \centering
  \caption{Quy mô các split NewsQA nguồn.}
  \label{tab:newsqa-source-splits}
  \begin{tabular}{@{}lrrr@{}}
    \toprule
    \textbf{Split} & \textbf{Số câu hỏi} & \textbf{Số context khác nhau}
    & \textbf{Câu hỏi/bài} \\
    \midrule
    Train & 74.160 & 10.895 & 6,8 \\
    Validation & 4.212 & 638 & 6,6 \\
    \bottomrule
  \end{tabular}
\end{table}

Từ 638 context của split validation, nhóm dùng seed 42 để chọn ngẫu nhiên 200
context và lấy toàn bộ câu hỏi tương ứng làm tập đánh giá. Split train có
10.895 context khác nhau; 31 context trùng nội dung với tập đánh giá bị loại để
tránh rò rỉ dữ liệu. 10.864 context còn lại được dùng làm distractor. Corpus
cuối gồm 11.064 bài báo: 200 bài đánh giá và 10.864 distractor. Các distractor
không có câu hỏi đánh giá, nhưng buộc hệ thống phải tìm bằng chứng trong một kho
dữ liệu lớn hơn.

\subsubsection{Hiệu chỉnh câu hỏi đánh giá}

Câu hỏi NewsQA được đặt khi người đọc đang xem bài báo, nên nhiều câu thiếu ngữ
cảnh khi đứng riêng. Nhóm giữ dữ liệu gốc để truy vết và tạo phiên bản
\texttt{resolved} theo quy trình sau:

\begin{enumerate}
  \item Codex đọc bài báo, gắn nhãn vấn đề và đề xuất sửa tối thiểu.
  \item Người đánh giá duyệt từng câu; phần thêm vào chỉ dùng để xác định chủ
  thể hoặc sự kiện, không chứa đáp án.
  \item Đáp án và evidence sai được sửa; câu không thể trả lời được loại.
  \item Các câu cùng bài, cùng ý nghĩa và cùng đáp án được gom cụm để loại trùng.
  \item Evidence span được ánh xạ sang \path{relevant_chunk_ids} sau khi chunk.
\end{enumerate}

Trong 1.340 câu nguồn, 4 câu không có đáp án đáng tin cậy bị loại. Sau review
còn 1.336 câu; loại tiếp 184 câu trùng nghĩa, thu được 1.152 câu
\texttt{resolved} thuộc 200 bài báo. Việc loại trùng diễn ra sau khi làm rõ câu
hỏi để so sánh đúng ý nghĩa đầy đủ của chúng.

\begin{table}[htbp]
  \centering
  \caption{Các mốc tạo tập câu hỏi đánh giá.}
  \label{tab:dataset-preparation}
  \begin{tabularx}{0.86\textwidth}{@{}Xrr@{}}
    \toprule
    \textbf{Mốc dữ liệu} & \textbf{Số câu} & \textbf{Thay đổi} \\
    \midrule
    Câu hỏi lấy từ 200 bài NewsQA & 1.340 & -- \\
    Sau human review & 1.336 & Loại 4 câu invalid \\
    Sau semantic deduplication & 1.152 & Loại 184 câu lặp \\
    \bottomrule
  \end{tabularx}
\end{table}

\begin{table}[htbp]
  \centering
  \caption{Các vấn đề được ghi nhận khi review.}
  \label{tab:question-review-issues}
  \begin{tabularx}{\textwidth}{@{}l r X@{}}
    \toprule
    \textbf{Issue code} & \textbf{Số lần} & \textbf{Ý nghĩa} \\
    \midrule
    \texttt{missing\_subject} & 688 & Thiếu chủ thể cần xác định \\
    \texttt{underspecified\_event} & 139 & Chưa xác định rõ sự kiện \\
    \texttt{generic\_reference} & 98 & Dùng tham chiếu chung như ``he'', ``it'' \\
    \texttt{wrong\_evidence} & 87 & Evidence không hỗ trợ đáp án \\
    \texttt{malformed\_question} & 71 & Câu hỏi sai hoặc thiếu cấu trúc \\
    \texttt{truncated\_answer} & 59 & Đáp án bị cắt thiếu \\
    \texttt{wrong\_answer} & 35 & Đáp án nguồn không đúng \\
    \texttt{yes\_no\_answer\_mismatch} & 4 & Câu hỏi Yes/No nhưng gold không phải Yes/No \\
    \bottomrule
  \end{tabularx}

  \vspace{1mm}
  {\footnotesize Một câu hỏi có thể mang nhiều issue code.}
\end{table}

Hai ví dụ về câu hỏi cần hiệu chỉnh:

\begin{enumerate}
  \item \textbf{Ngữ cảnh/làm rõ câu hỏi:} ``When was her surgery?'' được sửa
  thành ``When was Natalie Cole's kidney transplant surgery?''. Đáp án
  ``Tuesday'' không được đưa vào câu hỏi mới.
  \item \textbf{Đáp án:} câu hỏi về quân Canada có gold nguồn là ``35,000'',
  nhưng đây là tổng quân của các đồng minh NATO. Bài báo ghi Canada có hơn
  2.800 quân, nên nhóm sửa đáp án và evidence span, đồng thời giữ giá trị nguồn
  để truy vết.
\end{enumerate}

\subsection{Chia tập}

Toàn bộ partition được tạo theo bài báo. Danh sách article ID và question ID
được cố định bằng seed và SHA-256.

\begin{table}[htbp]
  \centering
  \caption{Các partition của tập resolved deduplicated.}
  \label{tab:evaluation-partitions}
  \begin{tabularx}{\textwidth}{@{}XrrX@{}}
    \toprule
    \textbf{Partition} & \textbf{Bài} & \textbf{Câu} & \textbf{Vai trò} \\
    \midrule
    Development & 50 & 281 & Screening và chọn cấu hình \\
    Held-out generation & 50 & 284 & Đánh giá cấu hình generation đã khóa \\
    Held-out reserve & 100 & 587 & Dữ liệu dự trữ, không dùng để tuning \\
    Phase 1 final-test retrieval & 150 & 871 & Toàn bộ held-out pool dùng kiểm tra retrieval \\
    \bottomrule
  \end{tabularx}
\end{table}

Trong 200 bài đánh giá, 50 bài với 281 câu được dùng làm development. 150 bài
còn lại tạo thành held-out pool gồm 871 câu. Từ 150 bài này, nhóm dùng seed 46
để chọn 50 bài với 284 câu cho held-out generation; 100 bài với 587 câu còn lại
được giữ làm reserve. Phase 1 dùng toàn bộ held-out pool 150 bài, còn Phase 2
chỉ dùng phần held-out generation 50 bài. Development và held-out pool không
trùng bài báo.

\subsection{Nguyên tắc thực nghiệm}

\subsubsection{Phase 1: đánh giá retrieval}

\begin{itemize}
  \item Đánh giá khả năng đưa gold evidence lên đầu danh sách; không gọi
  generator hoặc LLM-as-a-Judge.
  \item Các vòng chọn retriever, reranker và hybrid dùng cùng 281 câu
  development. Mỗi vòng chỉ thay đổi thành phần đang xét.
  \item Sau khi khóa cấu hình, final-test được chạy một lần trên 871 câu của 150
  bài chưa dùng để tuning. Kết quả này không được dùng để chọn lại cấu hình.
\end{itemize}

\subsubsection{Phase 2: đánh giá generation end-to-end}

\begin{itemize}
  \item Dùng cấu hình retrieval đã khóa từ Phase 1.
  \item Chạy retrieval và reranking một lần trên development, sau đó dùng lại
  cùng ranked traces cho mọi cấu hình generation.
  \item Khi thử context depth 1, 3 hoặc 5, pipeline chỉ lấy số context tương ứng
  từ đầu cùng một danh sách đã rerank.
\end{itemize}

Để kiểm soát chi phí API, Phase 2 được chia thành các chế độ sau:

\begin{table}[htbp]
  \centering
  \caption{Các chế độ chạy trong Phase 2.}
  \label{tab:phase2-run-modes}
  \begin{tabularx}{\textwidth}{@{}XrX@{}}
    \toprule
    \textbf{Chế độ} & \textbf{Số câu} & \textbf{Mục đích} \\
    \midrule
    Screening & 80 & So sánh nhanh prompt/context configuration \\
    Judge calibration & 20 & Tập con cố định của screening để chấm RAGAS \\
    Full development & 281 & Xác nhận finalist và áp dụng guardrail \\
    Held-out generation & 284 & Đánh giá cấu hình đã khóa trên 50 bài chưa thấy \\
    \bottomrule
  \end{tabularx}
\end{table}

\begin{itemize}
  \item Screening 80 câu và judge calibration 20 câu được chọn một lần với seed
  42 và giữ nguyên ID trong mọi thí nghiệm.
  \item Screening là bước tiết kiệm chi phí: loại cấu hình kém hơn trên tập nhỏ,
  chỉ giữ cấu hình có kết quả tương đương hoặc tốt hơn để làm finalist.
  \item Các finalist được chạy và chấm trên đủ 281 câu development. Guardrail
  được áp dụng ở bước này để chọn winner.
  \item Held-out chỉ được mở sau khi khóa winner và không dùng để tuning tiếp.
  \item Mỗi run lưu cấu hình, model, prompt hash, question IDs, source trace
  hash, coverage và retry log để có thể dừng và tiếp tục.
\end{itemize}

\subsection{Thước đo và kiểm định}

\subsubsection{Thước đo retrieval}

Với câu hỏi $q$, gọi $r_q$ là thứ hạng đầu tiên có gold chunk và $G_q$ là tập
gold chunks. Các thước đo được tính tại cutoff $k$:

\begin{align}
  \mathrm{Hit@}k(q) &= \mathbf{1}\!\left[G_q \cap R_q^{(k)} \neq \emptyset\right],\\
  \mathrm{Recall@}k(q) &= \frac{|G_q \cap R_q^{(k)}|}{|G_q|},\\
  \mathrm{MRR@}k(q) &=
  \begin{cases}
    1/r_q, & r_q \leq k,\\
    0, & r_q > k,
  \end{cases}
\end{align}

Trong đó, $R_q^{(k)}$ là $k$ kết quả đầu:

\begin{itemize}
  \item \textbf{Hit@$k$:} bằng 1 nếu top-$k$ chứa ít nhất một gold chunk; ngược
  lại bằng 0.
  \item \textbf{Recall@$k$:} tỷ lệ gold chunks xuất hiện trong top-$k$.
  \item \textbf{MRR@$k$:} dựa trên vị trí của gold chunk đầu tiên; gold xuất
  hiện càng sớm thì điểm càng cao.
  \item \textbf{nDCG@$k$:} đánh giá toàn bộ thứ tự top-$k$, ưu tiên gold chunks
  ở vị trí cao và chuẩn hóa theo thứ tự lý tưởng.
  \item \textbf{Latency:} thời gian truy xuất và rerank cho một câu hỏi.
  \item \textbf{Quy tắc chọn:} ưu tiên MRR@5; dùng nDCG@5, Hit@5 và latency để
  phân xử và mô tả trade-off.
\end{itemize}

\subsubsection{Thước đo generation và citation}

\begin{table}[H]
  \centering
  \caption{Vai trò của các thước đo Phase 2.}
  \label{tab:generation-metrics}
  \begin{tabularx}{\textwidth}{@{}p{42mm}X@{}}
    \toprule
    \textbf{Metric} & \textbf{Ý nghĩa} \\
    \midrule
    \multicolumn{2}{@{}l}{\textbf{So sánh đáp án}} \\
    \cmidrule(lr){1-2}
    Exact Match & Đáp án chuẩn hóa trùng hoàn toàn một accepted answer. \\
    Token F1 & Mức chồng lấp token giữa đáp án sinh và accepted answer tốt nhất. \\
    \addlinespace[2pt]
    \multicolumn{2}{@{}l}{\textbf{RAGAS}} \\
    \cmidrule(lr){1-2}
    Answer Correctness & RAGAS đánh giá độ đúng ngữ nghĩa so với gold answer. \\
    Answer Relevancy & Mức độ câu trả lời trực tiếp giải quyết câu hỏi. \\
    Faithfulness & Tỷ lệ các phát biểu trong đáp án được context hỗ trợ. \\
    Context Precision & Mức độ các context được đưa vào thực sự liên quan. \\
    Context Recall & Mức độ context bao phủ evidence cần để trả lời. \\
    \addlinespace[2pt]
    \multicolumn{2}{@{}l}{\textbf{Citation}} \\
    \cmidrule(lr){1-2}
    Citation Validity & Tỷ lệ citation trỏ tới một context index hợp lệ. \\
    Citation F1 & Trung bình điều hòa giữa precision và recall của citation so với gold evidence. \\
    \bottomrule
  \end{tabularx}
\end{table}

Exact Match, Token F1 và citation metrics được tính xác định từ artifact. Answer
Correctness, Answer Relevancy, Faithfulness, Context Precision và Context Recall
được chấm bằng RAGAS với judge tách biệt khỏi generator. Báo cáo phải ghi model,
provider, reasoning setting và score coverage của judge. Latency được báo cáo
theo P50/P95; token usage và chi phí chỉ so sánh khi các run dùng cùng phạm vi
câu hỏi và cùng cách tính giá.

\subsubsection{Khoảng tin cậy và điều kiện chọn cấu hình}

Khi so sánh hai cấu hình $A$ và $B$ trên cùng câu hỏi, nhóm tính sai khác theo
cặp:

\begin{equation}
  d_i = m(B,q_i)-m(A,q_i).
\end{equation}

\begin{itemize}
  \item \textbf{Paired bootstrap:} lấy lại các cặp sai khác $d_i$, thay vì lấy
  điểm của hai cấu hình độc lập.
  \item \textbf{Phase 1:} tạo 1.000 mẫu bootstrap; mỗi mẫu lấy lại 281 câu hỏi
  có hoàn lại.
  \item \textbf{Phase 2:} tạo 2.000 mẫu bootstrap; mỗi mẫu lấy lại 50 bài báo
  có hoàn lại và giữ toàn bộ câu hỏi của từng bài. Do đó, đây không phải 2.000
  tập gồm 50 câu hỏi.
  \item \textbf{CI95:} lấy percentile 2,5--97,5\% của 1.000 hoặc 2.000 giá trị
  trung bình mô phỏng.
  \item \textbf{Diễn giải:} CI95 không chứa 0 cho thấy khác biệt ở mức ý nghĩa
  0,05; CI95 chứa 0 nghĩa là chưa đủ bằng chứng về khác biệt, không phải hai cấu
  hình tương đương.
  \item \textbf{Guardrail:} kiểm tra delta trung bình so với baseline có nằm
  trong mức suy giảm cho phép hay không; đây là tiêu chí khác với CI95.
\end{itemize}

\begin{table}[htbp]
  \centering
  \caption{Guardrail dùng khi chọn cấu hình Phase 2 trên full development.}
  \label{tab:phase2-guardrails}
  \begin{tabular}{@{}lr@{}}
    \toprule
    \textbf{Điều kiện} & \textbf{Ngưỡng} \\
    \midrule
    $\Delta$ Faithfulness & $\geq -0{,}02$ \\
    $\Delta$ Citation F1 & $\geq -0{,}01$ \\
    $\Delta$ Citation Validity & $\geq -0{,}01$ \\
    Generation và judge coverage & $\geq 95\%$ \\
    \bottomrule
  \end{tabular}
\end{table}

Trong các ứng viên qua guardrail, Answer Correctness là metric chất lượng chính;
token cost, generation P95 và độ đơn giản của cấu hình lần lượt được dùng khi
khác biệt Answer Correctness nhỏ. Quy tắc này chỉ áp dụng trên full development.
Held-out được báo cáo nguyên trạng và không mở thêm vòng lựa chọn.
